\begin{lemma}
\label{editorial-source-lemma-radicial}
Let $\varphi : R \to S$ be a ring map. Assume
\begin{enumerate}
\item $\varphi$ induces an injective map of spectra,
\item $\varphi$ induces purely inseparable residue field extensions.
\end{enumerate}
Then for any ring map $R \to R'$ properties (1) and (2) are true for
$R' \to R' \otimes_R S$.
\end{lemma}

\begin{proof}
Retain the original base-change algebra
$S'=R'\otimes_R S$ and its two structural
maps. Let $\mathfrak p'\subset R'$ contract
to $\mathfrak p\subset R$, and set
$K=\kappa(\mathfrak p)$ and
$K'=\kappa(\mathfrak p')$ with the original
field embedding $K\to K'$. The original
fibre algebra is $F=K\otimes_R S$.
By Remark \ref{algebra-remark-fundamental-diagram},
its primes correspond to the primes
of $S$ over $\mathfrak p$, with their
original residue-field maps.

If there is no such prime, $F$ has
no prime and is zero, since every
nonzero unital ring has a maximal
ideal. The new fibre is zero as
well by the explicit comparison
below. Otherwise there is exactly
one original prime $\mathfrak q$
above $\mathfrak p$. Let $\mathfrak r$
be its corresponding prime in $F$.
It is the only prime of $F$, hence
is both its nilradical and its
unique maximal ideal. Consequently
$F/\mathfrak r$ is a field, and
the canonical map $F\to\kappa(\mathfrak r)$
is surjective with kernel $\mathfrak r$.
The original fibre correspondence
identifies this field with
$E=\kappa(\mathfrak q)$ over $K$.
By hypothesis $E/K$ is purely
inseparable. This proves the
needed nil-kernel quotient without
assuming that $F$ itself is reduced.

There is an original fibre isomorphism
$$
K'\otimes_{R'} S'\longrightarrow
K'\otimes_K F,\qquad
c\otimes(r'\otimes s)\longmapsto
c\overline r'\otimes(1\otimes s).
$$
Its inverse sends
$c\otimes(a\otimes s)$ to
$ca\otimes(1\otimes s)$, using the
specified embedding $K\to K'$.
Balancing follows from the original
commutative scalar diagram, and
these formulas compose to the
identity on every pure tensor.
Write $H=K'\otimes_K F$.
Lemma \ref{editorial-source-lemma-surjective-locally-nilpotent-kernel}
applied to $F\to E$ gives a
surjection
$H\to K'\otimes_K E$ with locally
nilpotent kernel.

Embed the purely inseparable $E$
in an algebraic closure of $K'$
by the unique roots compatible
with $K$. The full tensor comparison
proved in Lemma \ref{editorial-source-lemma-make-separable}
gives
$$
(K'\otimes_K E)_{red}\simeq C=K'E,
\qquad [\sum_i c_i\otimes e_i]
\longmapsto\sum_i c_i e_i.
$$
It applies also when $E/K$ is
infinite; each individual tensor
uses only finitely many root
exponents. The field $C/K'$ is
purely inseparable, because every
finite sum or quotient in the
compositum has a sufficiently high
$p$-power in $K'$ in positive
characteristic; in characteristic
zero $E=K$ and $C=K'$.
Both kernels in the composite
$H\to K'\otimes_K E\to C$ are nil.
Its full kernel is nil as well:
if an element maps to zero, some
power lies in the first nil kernel,
and a further power is zero.
Thus $H_{red}\simeq C$. The fibre
$H$ has exactly one prime and its
residue field is the original $C$.
The fibre correspondence proves
both required base-change properties.

In particular we obtain the exact
prime and residue field, rather
than only their uniqueness. Whenever
$\mathfrak q$ exists, the unique
$\mathfrak q'\subset S'$ over
$\mathfrak p'$ is the kernel of
$$
S'\longrightarrow C,\qquad
r'\otimes s\longmapsto
\overline r'\,\overline s,
$$
with the two factors embedded in
$K'$ and $E$ as above. Its contraction
to $S$ is $\mathfrak q$ and its
contraction to $R'$ is $\mathfrak p'$.
The fraction field of its image
is exactly $C$: it contains the
fractions giving $K'$ and $E$,
and is contained in their compositum.
This identifies $\kappa(\mathfrak q')$
with $C$ through the original maps.
It also proves that the image on
spectra after base change is exactly
the inverse image of the original
image on spectra.

Conversely, these two original
hypotheses are necessary for
injectivity on spectra after every
base change. Injectivity before
base change follows by choosing
$R'=R$. For an original
$\mathfrak q$ over $\mathfrak p$,
use the same $K,E,F$ above; injectivity
makes $F\to E$ a surjective nil-kernel
map. If $E/K$ were not purely
inseparable, the ring $E\otimes_K E$
would have at least two primes, as
we now prove with the original maps.
Its multiplication map onto $E$
has a prime kernel containing every
$z\otimes1-1\otimes z$.
If some $z\in E$ is transcendental
over $K$, no positive power of
this difference is zero: in its
full binomial expansion, the
coefficient of $1\otimes z^n$
is $(-1)^n$, and the powers
$1,z,\ldots,z^n$ remain independent
after extension from $K$ to $E$.

Otherwise $E/K$ is algebraic but
not purely inseparable. The
separable-first theorem supplies
$z\in E\setminus K$ separable over
$K$. Let $P(T)$ be its original
monic minimal polynomial, of degree
at least two. The ring
$E\otimes_K K(z)\simeq E[T]/(P)$
is reduced because $P$ is separable,
by the original polynomial and
Bezout argument in Lemma
\ref{editorial-source-lemma-separable-extension-preserves-reducedness}.
The element corresponding to $T-z$
is nonzero, by monic division,
and so is not nilpotent. The map
$E\otimes_K K(z)\to E\otimes_K E$
is injective by extension of a
vector-space basis. Thus the same
diagonal difference is not nilpotent
in the larger tensor ring either.
In both cases some prime avoids
that difference, while the original
multiplication kernel contains it.
They are two distinct primes.

Now take the actual base change
$R\to E$ through $R\to K\to E$.
Its algebra is
$E\otimes_R S\simeq E\otimes_K F$.
The surjection from this algebra
to $E\otimes_K E$ pulls the two
distinct primes back to distinct
primes, both over the unique prime
of the field $E$. This contradicts
universal injectivity. Hence every
original residue-field extension
must be purely inseparable, proving
the converse and the exact
characterization.
\end{proof}